%% Avant-propos
%%
%% L'avant-propos présente le contexte de publication des articles constitutifs
%% de la thèse, les contributions de chaque co-auteur et la pertinence des revues
%% choisies, conformément au Guide de présentation des mémoires et des thèses
%% (Polytechnique Montréal, janvier 2026) et aux exigences relatives aux thèses
%% par articles (section 3.4).

\ifthenelse{\equal{\Langue}{english}}{

\chapter*{FOREWORD}\thispagestyle{headings}
\addcontentsline{toc}{compteur}{FOREWORD}

}{
\chapter*{AVANT-PROPOS}\thispagestyle{headings}
\addcontentsline{toc}{compteur}{AVANT-PROPOS}

Cette thèse est présentée sous la forme d'une thèse par articles. Elle regroupe
trois articles de recherche originaux qui constituent le corps principal du
document (Chapitres~\ref{sec:Theme1}, \ref{sec:Theme2} et~\ref{sec:Theme3}).
Chaque article est reproduit intégralement dans sa version soumise ou acceptée,
reformaté pour assurer l'uniformité de la mise en page selon les exigences du
gabarit officiel de Polytechnique Montréal.

\section*{Article 1}

\noindent\textbf{Titre~:} \emph{Graph-Attentive Policy Fusion for Lightweight
Multi-Agent Routing in Disaster-Monitoring Wireless Sensor Networks (GAPF)}

\smallskip
\noindent\textbf{Auteur principal~:} Georges Parfait Djimefo Kapen

\smallskip
\noindent\textbf{Statut~:} Accepté pour publication après deux cycles de révision
en double-aveugle.

\smallskip
\noindent\textbf{Revue~:} \emph{IEEE Transactions on Green Communications and
Networking} (IEEE TGCN).\footnote{%
IEEE TGCN est une revue internationale à comité de lecture publiée par l'IEEE,
indexée notamment dans \textbf{Scopus} et le \textbf{Web of Science}, et
classée Q1 en «~Computer Networks and Communications~» selon Scimago.
Sa thématique couvre les communications vertes, l'efficacité énergétique et les
réseaux de nouvelle génération, en alignement direct avec les contributions de
cette thèse (routage écoénergétique, durabilité carbone, intégration 6G).
L’Article~1 est publié dans ce journal sous le DOI
\texttt{10.1109/TGCN.2026.3690507}.%
}

\smallskip
\noindent\textbf{Contribution de l'auteur~:} L'auteur est l'unique concepteur de
l'architecture GAPF (encodeur GCN, contrôleur CAS, réseau d'attention Q),
a réalisé l'ensemble des expériences et des analyses statistiques, et a rédigé
la totalité du manuscrit. Le directeur de recherche a fourni l'encadrement
scientifique et des commentaires sur les versions successives du manuscrit.

\section*{Article 2}

\noindent\textbf{Titre~:} \emph{Carbon-Aware Lifecycle-Autonomous Self-Healing
Routing for 6G-Integrated Disaster Sensor Networks: A Lifecycle-Sustainable
Approach (CALASH)}

\smallskip
\noindent\textbf{Auteur principal~:} Georges Parfait Djimefo Kapen

\smallskip
\noindent\textbf{Statut~:} Soumis à comité de lecture (en évaluation).

\smallskip
\noindent\textbf{Revue~:} \emph{IEEE Transactions on Green Communications and
Networking} (IEEE TGCN).

\smallskip
\noindent\textbf{Contribution de l'auteur~:} L'auteur est l'unique concepteur
des quatre piliers CALASH (CADR, CARE, SHDR, LSE) et de la métrique LCI,
a réalisé l'ensemble des simulations (30~graines Monte Carlo, 13~variantes
de protocoles), les analyses statistiques (Kruskal--Wallis, bootstrap) et
a rédigé la totalité du manuscrit. Le directeur de recherche a fourni
l'encadrement scientifique et des commentaires sur les versions successives
du manuscrit.

\section*{Article 3}

\noindent\textbf{Titre~:} \emph{Zero-Trust Shielded Graph-Structured Decision
Control for Secure Autonomous Recovery in AI-Native 6G Disaster-Monitoring
Sensing Networks (CASTER-ZT)}

\smallskip
\noindent\textbf{Auteur principal~:} Georges Parfait Djimefo Kapen

\smallskip
\noindent\textbf{Statut~:} Soumis à comité de lecture (en évaluation).

\smallskip
\noindent\textbf{Revue~:} \emph{IEEE Transactions on Mobile Computing} (IEEE TMC).

\smallskip
\noindent\textbf{Contribution de l'auteur~:} L'auteur est l'unique concepteur
du cadre CASTER-ZT (encodeur GraphSAGE, évaluateur double-hypothèse,
quantification d'incertitude MC-Dropout, bouclier déterministe gradué à
14~paramètres, dix propositions formelles), a réalisé la campagne de
2\,420~exécutions et rédigé la totalité du manuscrit. Le directeur de
recherche a fourni l'encadrement scientifique et des commentaires sur les
versions successives du manuscrit.

}